\documentclass{article}
\usepackage[UTF8]{ctex}
\usepackage{amsmath,amssymb,amsfonts}
\usepackage{enumitem}
\usepackage{geometry}
\geometry{a4paper, margin=1in}
% === 错题本排版优化（不参与题目标识） ===
\setlength{\parindent}{0pt}
\setlength{\parskip}{0.35em plus 0.10em minus 0.08em}
\setlength{\emergencystretch}{3em}
\setlist{itemsep=0.25em, topsep=0.40em, parsep=0pt}
\allowdisplaybreaks[2]
% === 错题本排版优化结束 ===
\begin{document}
\section*{安徽大学2026年高等代数真题}
\section*{105．安徽大学2026年高等代数试卷}
一．填空题．\\

\subsection{题目1}
1．设 \(f(x)=x^{3}+(t+1) x^{2}+2 x-4\) 与 \(g(x)=x^{3}+t x^{2}-2\) 的最大公因式为一个二次多项式，则 \(t=\)\\
\(\_\_\_\_\)．

\subsection{题目2}
2．设 \(A, B\) 都为可逆方阵，则 \(\left(\begin{array}{cc}A & C \\ O & B\end{array}\right)^{-1}=\) \(\_\_\_\_\)．

\subsection{题目3}
3．设 \(A\) 是三阶矩阵，已知 \(|A+I|=0,|A+2 I|=0,|A+3 I|=0\) ，则 \(|A+4 I|=\) \(\_\_\_\_\)．\\
\(\_\_\_\_\)

\subsection{题目4}
4．设 \(A\) 是 \(n\) 阶方阵，满足等式 \(A^{2}=2025 A+2026 I\) ，其中 \(I\) 为 \(n\) 阶单位矩阵。则有秩 \((A-2026 I)+\)秩 \((A+I)=\) \(\_\_\_\_\)。

\subsection{题目5}
5．已知列向量组 \(\left(\begin{array}{c}-3 \\ 1 \\ a\end{array}\right)\) 是方阵 \(\left(\begin{array}{ccc}-b & -2 & 1 \\ 2 & 3 & 1 \\ 3 & -1 & 4\end{array}\right)\) 的特征向量，则 \(4 a-3 b=\) \(\_\_\_\_\)．\\
\(\_\_\_\_\)

\subsection{题目6}
6．（可能有误）设矩阵 \(A\) 的初等因子为 \(\lambda, \lambda, \lambda^{3},(\lambda-1)^{2}, \lambda-2,(\lambda-2)^{2}\) ，且 \(A\) 的秩为 4 ，则 \(A\) 的所有不变因子为 \(\_\_\_\_\)．
\end{document}
